\honest{Final Words}{Eliezer Yudkowsky}{2009}

This is a story, the last in this book about the beisutsukai, the Bayesian Conspiracy. \nb{``Beisutsukai'' is a coined Japanese form of ``Bayes-user''.}

At dawn on Mount Mirror, Brennan and four other students wait for their teacher, Jeffreyssai, hoping for the ``real'' secrets. ``You're finished,'' he tells them. What he taught is incomplete, and they will harm themselves in trying to use it. They can reach mastery only by using what they have learned ``until they shatter in your hands.'' ``Go forth, then, and fail.''

Yin asks why he taught them at all. Brennan says that without teaching they would have no chance. Jeffreyssai agrees, then softens it: untaught, their chances would be ``less.'' A taught student who fails may recover and remake the Art, where an untaught one might think Reason itself had failed. ``The higher road begins after the Art seems to fail you; though the reality will be that it was you who failed your Art.'' \nb{Put as fact, this places every failure on the student and none on the Art, and the story gives no way to tell which is at fault in a given case. As advice on how to take a failure, it repeats ``Something to Protect.''} Is this the only road to mastery? ``I do not know,'' he says.

Then three flaws. The first is to look ``just the slightest bit harder for flaws in arguments whose conclusions you would rather not accept.'' The second is cleverness, which stays a weakness even when well known, because the clever value it for itself. The third is underconfidence that feels like humility, ending in lost momentum. \nb{The first is the thesis of ``Knowing About Biases Can Hurt People.'' The third is the subject of ``The Sin of Underconfidence,'' posted a week before this story, which calls it the third of three sins without naming the other two.} He leaves, saying: ``One of you at least seems likely to come back.''

Alone, Brennan finds that he wants nothing, and wonders whether his pursuit of power was only the role of ``an ambitious young man.'' The doubt he was taught has become the third flaw. Then he uses what he was taught: why would he not know a strong desire? Because he would not admit it, since it looks impossible. He calls this a ``silver-shoes moment.'' \nb{The phrase is from \textsc{The Wonderful Wizard of Oz}, where Dorothy learns that her silver shoes could have taken her home all along.} Pursuing his desire ``would make for a hard life, but not a sad one.'' I do not say what it is.

He sees that Jeffreyssai knew, and that what Taji said he wanted was weak evidence. He starts to think that ``it is futile to inquire how a beisutsukai master knows a thing,'' and stops: a would-be master should figure it out. ``It was, Brennan realized, a stupid proverb.'' \nb{The story has already shown the kind of evidence a master uses: Jeffreyssai read a flicker of Brennan's eyelids. The story's central claim, that mastery comes through failure, is foretold, not shown.}
